\documentclass[10pt,a4paper]{amsart}
\usepackage[margin=19mm]{geometry}
\usepackage{amsmath,amssymb,mathtools}
\usepackage[scr=rsfs,cal=euler]{mathalfa}
\usepackage{xfrac}
\usepackage[unicode,hidelinks,bookmarks=false]{hyperref}
\newcommand{\ie}{i.e.\ }
\newcommand{\cf}{cf.\ }
\newcommand{\resp}{respectively\ }
\newcommand{\Subsection}{\subsection}
\providecommand{\nobreakdash}{\nobreak\hbox{-}}
\setlength{\emergencystretch}{2em}
\multlinegap0pt
\usepackage{bm}
\usepackage{bbm}
\usepackage{dsfont}
\usepackage{xspace}
\usepackage{cancel}
\usepackage{mathtools}
\usepackage{amsthm}
\makeatletter
\@ifundefined{theorem}{\newtheorem{theorem}{Theorem}}{}
\@ifundefined{proposition}{\newtheorem{proposition}[theorem]{Proposition}}{}
\makeatother
\newcommand{\nperp}{\perp_{\delta}^\varepsilon}
\title[Norming Approximate Orthogonality in Normed Linear Spaces]{Norming Approximate Orthogonality in Normed Linear Spaces}
\author{Saikat Roy}
\date{Source: arXiv:2606.18727v2 (CC BY 4.0)}
\begin{document}
\maketitle

\noindent\emph{Source and licence.} Norming Approximate Orthogonality in Normed Linear Spaces. Version arXiv:2606.18727v2 (CC BY 4.0), distributed under the Creative Commons Attribution 4.0 International licence, as recorded in the arXiv licence metadata for this exact version and shown on the abstract page https://arxiv.org/abs/2606.18727v2. Full rights notice: licenses/arxiv-2606.18727-CC-BY-4.0.txt.\par\smallskip

\noindent\textit{Editorial scope.} Counted unit: the Proposition whose source label is \texttt{thm:bound} in the article (source0.tex lines 707--732, source label \texttt{thm:bound}), delivered together with the article's own complete original proof of it (source0.tex lines 734--802, one \texttt{proof} environment of 69 lines). No mathematics is written, repaired or re-derived: statement and proof are the article's own text line for line. Required context retained uncounted: Thm:Equiv (lines 182--201); thm:subspace (lines 222--233); prop:L (lines 293--304). Every definition, display and label of those lines is retained unchanged. Omitted from this package and not redistributed: 1--181, 202--221, 234--292, 305--706, 733--733, 803--1190. Nothing inside a retained excerpt is omitted or altered: every retained body is delivered exactly as the article's lines are written, so no omission receipt of any kind accompanies this package. Citations: the retained excerpts carry no citation command at all, so no bibliography entry of the article is transcribed and no reference is redistributed. Reference adaptation: none was needed and none was made; every reference inside the delivered unit resolves inside it, so no unresolved reference or citation is produced. Printed slips kept literal: no printed word, symbol, hypothesis or number of the counted statement or of its proof was altered. Layout and class: the article's own class is \texttt{amsart} with packages including mathtools, amssymb, inputenc, tikz, pgfplots, mathrsfs, mathptmx, xcolor, hyperref; the wrapper is the lane's pinned \texttt{amsart} editorial wrapper at 10pt on A4, which restates the theorem-like environments the retained text needs and adds the article's own macro definitions that the retained text needs (\texttt{\textbackslash nperp}), with extra wrapper packages (bm, bbm, dsfont, xspace, cancel, mathtools). The delivered numbering is the unit's own. Assets: no retained span contains a figure environment, a graphics-inclusion command, a PDF-page inclusion, a table or a TikZ picture; no image file is copied into this package, the package declares no asset dependency and no image file is redistributed with it. Licence: version arXiv:2606.18727v2 of this article is distributed under the Creative Commons Attribution 4.0 International licence, as recorded in the arXiv licence metadata for this exact version and shown on the abstract page https://arxiv.org/abs/2606.18727v2. A full rights notice is delivered at licenses/arxiv-2606.18727-CC-BY-4.0.txt. Characters: every retained excerpt is pure ASCII, so the delivered engine, XeLaTeX with its default Latin Modern text font, reports no missing character.
\begin{theorem}\label{Thm:Equiv}
Let $X$ be a normed linear space, $x, y \in X$ non-zero, and $\delta, \varepsilon 
\in [0,1)$ with $\varepsilon < (1-\delta)^2$. The following conditions are equivalent:
\begin{itemize}
    \item[(i)] $x \nperp y$.
    \item[(ii)] $\exists\, f \in X^*$ with $f(x)$ real and $f(x) \geq (1-\delta)\|f\|\|x\|$ 
    and $|f(y)| \leq \frac{\varepsilon}{1-\delta}\|f\|\|y\|$.
    \item[(iii)] $\exists\, f \in X^*$ with $\mathrm{Re}\,f(x) \geq (1-\delta)\|f\|\|x\|$ 
    and $|f(y)| \leq \frac{\varepsilon}{1-\delta}\|f\|\|y\|$.
    \item[(iv)] $\|x + \lambda y\| \geq (1-\delta)\|x\| - 
    \frac{\varepsilon}{1-\delta}|\lambda|\|y\|$ for all scalars $\lambda$.
\end{itemize}
Moreover, if $X$ is a Hilbert space, then \emph{(i)--(iv)} are further equivalent to:
\begin{itemize}
    \item[(v)] $|\langle x, y\rangle| \leq L(\delta,\varepsilon)\|x\|\|y\|$, where
$L(\delta,\varepsilon)=\cos(\theta_\varepsilon-\theta_\delta)$ is the constant of Proposition~\ref{prop:L}.
   \item[(vi)] $\exists\, u \in S_X$ with $\langle x, u\rangle \geq (1-\delta)\|x\|$ 
    and $|\langle y, u\rangle| \leq \frac{\varepsilon}{1-\delta}\|y\|$.
\end{itemize}
\end{theorem}

\begin{theorem}\label{thm:subspace}
Let $Y$ be a closed subspace of $X$ and $x \in X \setminus Y$. Let $\varepsilon, \delta \in [0,1)$ be such that $\varepsilon < (1-\delta)^2$. Then 
the following are equivalent:
\begin{itemize}
    \item[(a)] $x \nperp y$ 
for every $y \in Y$;
    \item[(b)]  there exists $f \in X^*$ with $\|f\| = 1$ such that
\[
|f(x)| \geq (1-\delta)\|x\|, \qquad \|f|_Y\| \leq \frac{\varepsilon}{1-\delta}.
\]
\end{itemize}
\end{theorem}

\begin{proposition}\label{prop:L}
Let $\delta,\varepsilon \in [0,1)$ with $\varepsilon < (1-\delta)^2$, and define
\[
L(\delta,\varepsilon) := \cos(\theta_\varepsilon-\theta_\delta).
\]
Then
\[
L(\delta,\varepsilon) = \varepsilon +
\frac{\sqrt{\delta(2-\delta)\bigl((1-\delta)^2-\varepsilon^2\bigr)}}{1-\delta}
\qquad\text{and}\qquad L(\delta,\varepsilon)<1 .
\]
\end{proposition}

\begin{proposition}\label{thm:bound}
Let $Y$ be a closed proper non-trivial subspace of $X$, $x \in X \setminus \{0\}$,
and $\delta \in [0,1)$. Put
\[
\widetilde m_\delta(x,Y) := \inf\bigl\{\|f|_Y\| : f \in S_{X^*},\
|f(x)| \geq (1-\delta)\|x\|\bigr\},
\]
\[
m_\delta(x,Y) := \inf\bigl\{\|f|_Y\| : f \in B_{X^*},\
|f(x)| \geq (1-\delta)\|x\|\bigr\},
\]
and suppose $\widetilde m_\delta(x,Y) < 1-\delta$. Then we have
\begin{itemize}
    \item[(a)] $x\notin Y$, both infima are attained.
    \item[(b)] The set
\[
S := \bigl\{\varepsilon \in [0,(1-\delta)^2) : x \perp_\delta^{\varepsilon} Y\bigr\}
\]
has a least element $\varepsilon_{\min}$, and
\[
\varepsilon_{\min} = (1-\delta)\,\widetilde m_\delta(x,Y).
\]
    \item[(c)] $(1-\delta)\,m_\delta(x,Y) \leq \varepsilon_{\min} \leq m_\delta(x,Y)$, and
$x \perp_\delta^{m_\delta(x,Y)} Y$ whenever $m_\delta(x,Y)<(1-\delta)^2$.
\end{itemize}
\end{proposition}

\begin{proof}
Define $\mathscr{N}_\delta(x) := \{ f \in B_{X^*} : |f(x)| \geq (1-\delta)\|x\| \}$.
This set is a weak$^*$-closed subset of $B_{X^*}$, hence weak$^*$-compact by the
Banach--Alaoglu theorem and it is non-empty by the Hahn--Banach
theorem. Since the map $f \mapsto \|f|_Y\|$ is
weak$^*$-lower semicontinuous, it attains its infimum on $\mathscr{N}_\delta(x)$,
so $m_\delta(x,Y)$ is well-defined and the minimum is attained.


\medskip

\noindent The hypothesis forces $x\notin Y$. Indeed, if $x\in Y$, then every $f\in S_{X^*}$ with $|f(x)|\geq(1-\delta)\|x\|$ satisfies $\|f|_Y\|\,\|x\|\geq |f(x)|\geq (1-\delta)\|x\|$, proving $\widetilde m_\delta(x,Y)\geq 1-\delta$, a contradiction.


\medskip


\noindent Let $\varepsilon_* := (1-\delta)\,\widetilde m_\delta(x,Y)$,
so that $\varepsilon_*<(1-\delta)^2$ by hypothesis. Since $x\notin Y$,
Theorem~\ref{thm:subspace} gives, for every $\varepsilon\in[0,(1-\delta)^2)$,
\[
x \perp_\delta^{\varepsilon} Y \iff \exists\, f\in S_{X^*}:~ |f(x)|\geq (1-\delta)\|x\|,~
\|f|_Y\|\leq \tfrac{\varepsilon}{1-\delta}.
\]
If $\varepsilon>\varepsilon_*$ then $\widetilde m_\delta(x,Y)<\varepsilon/(1-\delta)$, so
such an $f$ exists and $\varepsilon\in S$; conversely, if $\varepsilon\in S$ then by Theorem \ref{thm:subspace}, $\widetilde m_\delta(x,Y)\leq \varepsilon/(1-\delta)$, that is
$\varepsilon\geq\varepsilon_*$. Hence $S\neq\emptyset$ and $\inf S=\varepsilon_*$.


\medskip

\noindent To see that the infimum is attained, take any sequence
$\varepsilon_n \searrow \varepsilon_*$
with each $\varepsilon_n \in S$. By the equivalence
$(i)\Leftrightarrow(iv)$ of Theorem~\ref{Thm:Equiv}, for every scalar $\lambda$
and every $y \in Y$,
\[
\|x + \lambda y\|
\geq (1-\delta)\|x\| - \frac{\varepsilon_n}{1-\delta}|\lambda|\|y\|
\xrightarrow{n \to \infty}
(1-\delta)\|x\| - \frac{\varepsilon_*}{1-\delta}|\lambda|\|y\|,
\]
so condition $(iv)$ holds at $\varepsilon_*$ for every $y \in Y$, and by
$(iv)\Rightarrow(i)$ of Theorem~\ref{Thm:Equiv} we conclude
$x \perp_\delta^{\varepsilon_*} Y$. Therefore, $\varepsilon_{\min}=\varepsilon_*=(1-\delta)\widetilde m_\delta(x,Y)$, and there exists $h\in S_{X^*}$ with
$|h(x)|\geq(1-\delta)\|x\|$ and $\|h|_Y\|\leq \varepsilon_{\min}/(1-\delta)=\widetilde m_\delta(x,Y)$. This also shows
\[
\widetilde m_\delta(x,Y) := \min\bigl\{\|f|_Y\| : f \in S_{X^*},\
|f(x)| \geq (1-\delta)\|x\|\bigr\},
\]

\medskip

\noindent Evidently, $m_\delta(x,Y)\leq \widetilde m_\delta(x,Y)$, since $S_{X^*}\subset B_{X^*}$. Conversely, let $f_0 \in \mathscr{N}_\delta(x)$ attain $m_\delta(x,Y)$. From
$\|f_0\|\,\|x\|\geq|f_0(x)| \geq (1-\delta)\|x\|$ we get $\|f_0\| \geq 1-\delta > 0$, so
$g_0 := f_0/\|f_0\| \in S_{X^*}$ satisfies
\[
|g_0(x)| = \frac{|f_0(x)|}{\|f_0\|} \geq (1-\delta)\|x\|,\qquad
\|g_0|_Y\| = \frac{\|f_0|_Y\|}{\|f_0\|}\leq \frac{m_\delta(x,Y)}{1-\delta},
\]
proving $\widetilde m_\delta(x,Y)\leq m_\delta(x,Y)/(1-\delta)$. Multiplying the chain
$m_\delta(x,Y)\leq \widetilde m_\delta(x,Y)\leq m_\delta(x,Y)/(1-\delta)$ by $(1-\delta)$
gives
\[
(1-\delta)\,m_\delta(x,Y) \leq \varepsilon_{\min} \leq m_\delta(x,Y).
\]
If in addition $m_\delta(x,Y)<(1-\delta)^2$, then $m_\delta(x,Y)\in S$ by monotonicity, so
$x \perp_\delta^{m_\delta(x,Y)} Y$. This completes the proof.
\end{proof}

\end{document}
